% ===================================================================
%  Αρχείο: 13537_SOLUTION.tex (Fragment)
%  Αυτόματη μετατροπή μέσω doc2latex.py
% ===================================================================

ΘΕΜΑ 4

Στο παρακάτω σχήμα δίνεται ισοσκελές τρίγωνο ΑΒΓ με $ΑΒ = Α\Gamma\ $, σημείο Δ της πλευράς ΑΓ, ώστε $Α\Delta = Β\Delta = Β\Gamma$ και σημείο Ε της πλευράς ΑΒ, ώστε $ΑΕ = \Gamma\Delta$.

\begin{alist}
\item Να αποδείξετε ότι:

.}
\end{alist}
\begin{alist}
\item
  $\widehat{\Gamma} = 2\widehat{Α}$ . \monades{6}
\item
  $\widehat{Α} = 36^{ο}$. \monades{6}
\item
  Το τρίγωνο ΑΔΕ είναι ισοσκελές. \monades{6}

\item Στην προέκταση της ΔΕ προς το Ε θεωρούμε σημείο Ζ, ώστε $\Delta Ζ = Α\Gamma$. Να αποδείξετε ότι το τρίγωνο ΒΔΖ είναι ισοσκελές. \monades{7}


% TODO: Μετατροπή σχήματος σε TikZ/PGFPlots
\begin{figure}[htbp]
\centering
\includegraphics[width=0.8\linewidth]{../extracted/13537_SOLUTION/image1.pdf}
% \caption{Σχήμα 1}
\end{figure}
\end{alist}
